\section{Software Development Processes}
\subsection{Part I: Introduction}
\subsubsection*{Warum brauchen wir einen Software-Engineering-Prozess?}

    \begin{enumerate}
        \item \fatsf{Größe der Aufgabe}
        \begin{itemize}
            \item Organisation eines potenziell großen Teams
            \item Zuweisung von Verantwortlichkeiten
            \item Definition von Kollaborationsmethoden
        \end{itemize}
        \item \fatsf{Komplexität der Aufgabe}
        \begin{itemize}
            \item Viele verschiedene Aktivitäten
            \item Abhängigkeiten zwischen den Aufgaben (Wann wird was erledigt?)
        \end{itemize}
        \item \fatsf{Qualitätskontrolle}
        \begin{itemize}
            \item Sicherstellen, ob das System korrekt funktioniert
            \item Erkennen und Beheben von Fehlern
        \end{itemize}
    \end{enumerate}

\subsubsection*{Was ist der Software-Engineering-Prozess?}

    Ein \fatsf{Software-Enginnering-Prozess} ist die Menge an Aktivitäten und Ergebnissen, die zur Erstellung eines Softwareprodukts führen\\
    \fatsf{Grundlegende Prozessaktivitäten:}
    \begin{enumerate}
        \item \fatsf{Software-Spezifikation}
        \begin{itemize}
            \item Definition der zu entwickelnden Software
            \item Festlegung der betrieblichen Rahmenbedingungen
        \end{itemize}
        \item \fatsf{Software-Entwicklung}
        \begin{itemize}
            \item Design, Implementierung und Verifikation des Codes
        \end{itemize}
        \item \fatsf{Software-Validierung}
        \begin{itemize}
            \item Sicherstellen, dass die Software den Kundenanforderungen entspricht
        \end{itemize}
        \item \fatsf{Software-Evolution}
        \begin{itemize}
            \item Anpassung und Modifikation an neue Kunden- oder Marktanforderungen
        \end{itemize}
    \end{enumerate}

\subsubsection*{Software-Engineering-Prozessmodelle}

    Software-Prozessmodelle sind \fatsf{abstrakte Beschreibungen}, die eine bestimmte Sicht auf den Entwicklungsprozess bieten\\
    Laut lan Sommerville beinhalten Prozessmodelle:
    \begin{itemize}
        \item Entwicklungsaktivitäten
        \item Softwareprodukte (z.B. Architektur, Code, Dokumentation)
        \item Rollen der beteiligten Personen
    \end{itemize}
    \fatsf{Beispiele für Software-Prozessmodelle:}
    \begin{itemize}
        \item \fatsf{Wasserfallmodell} (linear, sequenzieller Entwicklungsprozess)
        \item \fatsf{Spiralmodell} (iterativer, risikogesteuerter Ansatz)
        \item \fatsf{V-Modell (XT)} (erweiterte Version des Wasserfallmodells mit Verifikations- und Validierungsschritten)
        \item \fatsf{Extreme Programming (XP)} (agile Modell mit Fokus auf schnelles Feedback und kontinuierliche Verbesserung)
    \end{itemize}
    \fatsf{Hinweis:} Große Projekte verwenden oft eine \fatsf{Kombination mehrerer Prozessmodelle}, um unterschiedliche Softwarekomponenten zu entwickeln

\subsection{Part II: The Waterfall Model}
\subsubsection*{Das Wasserfallmodell}
\begin{figure}[h]
    \centering
    \includegraphics[width=0.5\textwidth]{Software Development Processes Images/The Waterfall model.jpeg}
\end{figure}

    Ein \fatsf{generischer Prozessmodell}, das in fünf aufeinanderfolgende Phasen unterteilt ist:
    \begin{enumerate}
        \item \fatsf{Anforderungsanalyse und -definition}
        \begin{itemize}
            \item Anforderungen werden mit den Nutzern abgestimmt
            \item Detaillierte Spezifikation der Systemanforderungen
        \end{itemize}
        \item \fatsf{System- und Softwaredesign}
        \begin{itemize}
            \item Definition der Systemarchitektur
            \item Identifikation der grundlegenden Abstraktionen und ihrer Beziehungen
        \end{itemize}
        \item \fatsf{Implementierung und Unit-Testing}
        \begin{itemize}
            \item Umsetzung des Designs in Programmcode
            \item Tests überprüfen, ob jede Einheit der Spezifikation entspricht
        \end{itemize}
        \item \fatsf{Integration und Systemtest}
        \begin{itemize}
            \item Einzelne Programmteile werden zusammengefügt und als Gesamtsystem getestet
        \end{itemize}
        \item \fatsf{Betrieb und Wartung}
        \begin{itemize}
            \item Nach dem Release wird das System gewartet und weiterentwickelt
        \end{itemize}
    \end{enumerate}
    \fatsf{Merkmale des Wasserfallmodells:}
    \begin{itemize}
        \item Jede Phase liefert geprüfte Artefakte (Dokumente, Code, Spezifikationen)
        \item Die nächste Phase beginnt erst nach Abschluss der vorherigen
        \item Fehler führen dazu, dass vorherige Phasen wiederholt werden müssen
        \item Ähnelt traditionellen Ingenieursprozessen (z.B. Maschinenbau)
    \end{itemize}

\subsubsection*{Kritik am Wasserfallmodell}

    \begin{itemize}
        \item \fatsf{Nicht iterativ} $\rightarrow$ Änderungen in späten Phasen sind schwierig und teuer
        \item \fatsf{Frühe Prototypen nicht möglich} $\rightarrow$ Probleme erst spät erkannt
        \item \fatsf{Testen erst am Ende} $\rightarrow$ Fehler früh im Prozess bleiben unentdeckt
        \item \fatsf{Unterschiedliche Teams in verschiedenen Phasen} $\rightarrow$ Wissen geht verloren
    \end{itemize}
    \fatsf{Zitat:} Bereits 1970 als problematisches Modell für Softwareentwicklung erkannt (Winston Royce, Managing the Development of Large Software Systems)

\subsubsection*{Das V-Modell}

    Eine verbesserte Version des Wasserfallmodells mit stärkerem Fokus auf Tests.\\
    Jede Entwicklungsphase hat eine zugehörige Testphase:
    \begin{itemize}
        \item Anforderungsspezifikation $\rightarrow$ Akzeptanztestplan
        \item Systemdesign $\rightarrow$ Systemintegrationtestplan
        \item Detailliertes Design $\rightarrow$ Subsystem-Integrationstestplan
        \item Implementierung $\rightarrow$ Modul-/Unit-Testplan
    \end{itemize}
    \fatsf{Merkmale:}
    \begin{itemize}
        \item Klare Struktur
        \item Strenge Dokumentation
        \item Gut geeignet für sicherheitskritische Systeme (z.B. Medizintechnik, Luftfahrt)
    \end{itemize}

\subsubsection*{Kritik am V-Modell}

    \begin{itemize}
        \item Nicht iterativ $\rightarrow$ Änderungen bleiben teuer
        \item Unflexibel bei sich ändernden Anforderungen
        \item Nur geeignet für kleine, stabile Projekte mit klaren Anfroderungen
    \end{itemize}

\subsection{Part III: Agile Development}
\subsubsection*{Agile Softwareentwicklung}

    \fatsf{Kernpunkte (Key Points)}
    \begin{itemize}
        \item Wartungszentrierter Prozess
        \item Schnelle Entwicklung trotz sich ändernder Anforderungen
        \item Kurze und schnelle Entwicklungszyklen (\enquote{agil})
        \item Ursprünglich für kleine bis mittelgroße Teams (3-9 Personen) entwickelt
    \end{itemize}
    \fatsf{Grundannahmen (Assumptions)}
    \begin{itemize}
        \item \fatsf{Enge Zusammenarbeit} zwischen Analysten, Kunden, Entwicklern und Testern
        \item \fatsf{Disziplinierte Praktiken und kontinuierliches Feedback} notwendig
        \item \fatsf{Flexibles und wartbares Design} durch bewährte Designprinzipien
        \item \fatsf{Agilität ersetzt nicht die Validierung!} Der Kunde muss als Teammitglied mitwirken
    \end{itemize}

\subsubsection*{Agiles Manifest (Manifesto for Agile Software Development)}
\begin{enumerate}
    \item \fatsf{Individuen \& Interaktionen statt Prozesse \& Tools}
    \begin{itemize}
        \item Das beste Tool nützt nichts, wenn das Team nicht zusammenarbeitet
        \item Kleine Teams $\rightarrow$ wachsen \& schrumpfen je nach Bedarf
        \item Regelmäßige Reflexionen über Prozesse \& Arbeitsumfeld
    \end{itemize}
    \item \fatsf{Funktionierende Software statt umfassende Dokumentation}
    \begin{itemize}
        \item \fatsf{Dokumentation soll inkrementell erfolgen, nicht upfront}
        \item Kein Vorwand, um Doku ganz zu überspringen
        \item \fatsf{Code spielt zentrale Rolle} und muss sagar besser dokumentiert sein
    \end{itemize}
    \item \fatsf{Zusammenarbeit mit dem Kunden statt Vertragsverhandlungen}
    \begin{itemize}
        \item Verträge sollen \fatsf{Zusammenarbeit mit dem Kunden definieren}, nicht nur feste Fristen und Preise
        \item \fatsf{Feste Preis- \& Terminverträge scheitern oft}, da Anforderungen sich ändern
        \item Der Kunde will aber trotzdem wissen, was er für sein Geld bekommt
    \end{itemize}
    \item \fatsf{Reaktion auf Veränderungen statt striktes Befolgen eines Plans}
    \begin{itemize}
        \item \fatsf{Pläne sind wichtig, aber müssen flexiblel bleiben}
        \item Anforderungen ändern sich $\rightarrow$ Prozesse müssen anpassbar sein
    \end{itemize}
\end{enumerate}
\subsubsection*{Agile Methoden \& Prozesse}
Es gibt viele agile Methoden, darunter:
\begin{itemize}
    \item Extreme Programming (XP)
    \item SCRUM (Projektmanagement-Methode)
    \item Agile Unified Process (AUP)
    \item Feature-Driven Development (FDD)
    \item Crystal Methodology
    \item Adaptive Software Development (ASD)
\end{itemize}
\begin{figure}[h]
    \centering
    \includegraphics[width=1\textwidth]{Software Development Processes Images/Scope und Level of Support.jpeg}
\end{figure}
\subsection{Part IV: Extreme Programming (XP)}

    XP besteht aus einer Reihe von \fatsf{einfachen, aber stark vernetzten Praktiken}, die darauf abzielen, Software effizient und flexibel zu entwickeln

\subsubsection*{Kunde als Teil des Teams}
\begin{itemize}
    \item \fatsf{Kunde definiert \& priorisiert Features}
    \item Ist \fatsf{ständig verfügbar} für Rückfragen
    \item Falls nicht möglich $\rightarrow$ Team kann einen \fatsf{Produktverantwortlichen (Product Owner)} benennen
\end{itemize}
\subsubsection*{User Stories}
\begin{itemize}
    \item \fatsf{Kurzbeschreibung} der Anforderungen durch den Kunden
    \item Enthält eine \fatsf{Schätzgröße (Strory Points)} zur Aufwandseinschätzung
    \item \fatsf{Keine technischen Datails}, da sich diese oft ändern
\end{itemize}
\fatsf{Gute User Stories folgen INVEST:}
\begin{itemize}
    \item Independent $\rightarrow$ Unabhängig von anderen Stories
    \item Negotiable $\rightarrow$ Bis zum Sprintbeginn verhandelbar
    \item Valuable $\rightarrow$ Brint einen echten Nutzen
    \item Estimable $\rightarrow$ Aufwand kann eingeschätzt werden
    \item Small $\rightarrow$ Klein genug, um pro Iteration mehrere zu implementieren
    \item Testable $\rightarrow$ Muss überprüfbar sein
\end{itemize}
\fatsf{Beispiel-Template:} \enquote{Als [Rolle] möchte ich [Ziel], damit [Nutzen]}
\subsubsection*{Akzeptanztests (Acceptance Tests)}
\begin{itemize}
    \item Tests werden \fatsf{vor oder während der Implementierung} erstellt
    \item Sobald ein Akzeptanztest bestanden wird, \fatsf{darf er nie wieder fehlschlagen} (Regression vermeiden)
\end{itemize}
\subsubsection*{Kurze Entwicklungszyklen (Short Development Cycles)}
\begin{itemize}
    \item Iterationen sind \fatsf{zeitlich festgelegt} (time-boxed)
    \item Am Ende jeder Iteration gibt es eine \fatsf{testbare Software}
    \item Falls Features nicht rechtzeitig festig werden $\rightarrow$ \fatsf{verschieben auf spätere Iterationen}
\end{itemize}
\subsubsection*{Das Planungsspiel (The Planning Game)}
\begin{itemize}
    \item \fatsf{Geschäftsseite (Customer):} Bestimmt \fatsf{Wichtigkeit \& Priorität}
    \item \fatsf{Entwicklerseite:} Bestimmt \fatsf{technischen Aufwand \& Konsequenzen}
    \item \fatsf{Ergebnis:}
    \begin{itemize}
        \item Kunde wählt Features für den Release
        \item Entwickler schätzen Aufwand \& setzen Zeitpläne fest
    \end{itemize}
\end{itemize}
\fatsf{Iterative Planung}
\begin{enumerate}
    \item \fatsf{Initial Exploration (Erste Erkundung)}
    \begin{itemize}
        \item \fatsf{Ziel:} Erste Sammlung der wichtigsten \fatsf{User Stories}
        \item \fatsf{Wie?}
        \begin{itemize}
            \item Entwickler und Kunden identifizieren \fatsf{die wichtigsten} Anforderungen in Form von \fatsf{User Stories}
            \item Entwickler schätzen die User Stories mit \fatsf{Story Points}
            \item Story Points messen den \fatsf{relativen Aufwand}, nicht die absolute Zeit
            \item Beispiel: Eine Story mit \fatsf{4 Punkten} ist doppelt so aufwändig wie eine mit \fatsf{2 Punkten}
        \end{itemize}
        \item \fatsf{Besonderheit:} Da am Anfang noch keine Erfahrungswerte existieren, basiert die Schätzung \fatsf{zunächst auf Erfahurng und Annahmen}
        \item Später wird die Schätzung mit der \fatsf{tatsächlichen Geschwindigkeit (Velocity)} verglichen und verbessert
    \end{itemize}
    \item \fatsf{Release-Planung:}
    \begin{itemize}
        \item \fatsf{Ziel:} Festlegung des \fatsf{ersten Release-Termins} in \fatsf{2-4 Monaten}
        \item \fatsf{Wie?}
        \begin{itemize}
            \item Kunde wählt User Stories aus, die bis zum ersten Release fertig sein sollen
            \item Die Anzahl der User Stories wird durch die \fatsf{aktuellen Erfahrungswerte zur Velocity} bestimmt
            \item \fatsf{Velocity} = Anzahl der in einer Iteration \fatsf{abgeschlossenen Story Points}
        \end{itemize}
        \item \fatsf{Wichtig:}
        \begin{itemize}
            \item Da die Velocity in den ersten Iterationen noch nicht stabil ist, wird der Plan später \fatsf{angepasst}
        \end{itemize}
    \end{itemize}
    \item \fatsf{Iterations-Planung:}
    \begin{itemize}
        \item \fatsf{Ziel:} Auswahl der User Stories für die nächste Iteration
        \item \fatsf{Wie?}
        \begin{itemize}
            \item Kunde wählt User Stories für die nächste Iteration (Sprint)
            \item \fatsf{Die Summer der Story Points darf die Velocity der vorherigen Iteration nicht übersteigen}
            \item Innerhalb der Iteration \fatsf{dürfen Entwickler die Reihenfolge selbst bestimmen}
            \item Die Iteration endet \fatsf{fix nach einer bestimmten Zeit (z.B. 2 Wochen)}, selbst wenn nicht alle geplanten Stories fertig sind
        \end{itemize}
    \end{itemize}
    \item \fatsf{Task Planning (Deatilplanung der Iteration)}
    \begin{itemize}
        \item \fatsf{Ziel:} Zerlegen der User Stories in kleinere \fatsf{Tasks}
        \item \fatsf{Wie?}
        \begin{itemize}
            \item Jede User Story wird in einzelne \fatsf{Entwicklungsaufgaben (Tasks)} aufgteilt
            \item \fatsf{Dauer einer Aufgabe: 4-16 Stunden}
            \item Entwickler \fatsf{wählen frei} aus, welche Aufgaben sie übernehmen
            \item Auch Nicht-Experten können Aufgaben übernehmen $\rightarrow$ \fatsf{Lernen wird gefördert}
        \end{itemize}
    \end{itemize}
\end{enumerate}
\subsubsection*{Berechnung der Velocity in XP}
Die Velocity kann auf zwei Arten berechnet werden:

\fatsf{1. Standardmethode}
\[
\text{Velocity} = \frac{\text{Story Points der abgeschlossenen User Stories}}{\text{Anzahl der Iterationen}}
\]

\fatsf{2. Alternative Methode}
\[
\text{Velocity} = \frac{\text{Story Points der abgeschlossenen User Stories}}{\text{Gesamtzeit (in Stunden)}}
\]

\fatsf{Beispiel}
In Iteration 1 wurden zwei User Stories abgeschlossen:
\begin{itemize}
    \item „Verfügbare Autos anzeigen“ (2 Story Points, 1,5h Aufwand)
    \item „Auto buchen“ (2 Story Points, 3,5h Aufwand)
\end{itemize}

Der gesamte Zeitaufwand beträgt:
\[
1,5h + 3,5h = 5h
\]

Die alternative Velocity-Berechnung ergibt:
\[
\frac{2 + 2}{1,5 + 3,5} = 0,8 \text{ Story Points pro Stunde}
\]

\fatsf{Fazit}
\begin{itemize}
    \item Die Velocity gibt an, \textbf{wie viele Story Points pro Iteration} umgesetzt werden können.
    \item Sie wird mit der Zeit immer genauer.
    \item Sie hilft bei der \textbf{Planung zukünftiger Iterationen und Releases}.
\end{itemize}
\subsubsection*{Simples Design (Simple Design)}
\begin{itemize}
    \item Fokus auf \fatsf{das aktuell Notwendige} $\rightarrow$ Kein Over-Engineering!
    \item \enquote{\fatsf{You aren't gonna need it!} (YAGNI)}
    \item Kein redundanter Code $\rightarrow$ \fatsf{Nutzung von Abstraktionen \& Design Patterns}
\end{itemize}
\subsubsection*{Testgetriebene Entwicklung (Test-Driven Development, TDD)}
\begin{itemize}
    \item \fatsf{Schreib zuerst Tests, dann Code!}
    \item Ziel: Jede Funktionalität wird erst implementiert, wenn ein Test sie überprüft
    \item Vorteile:
    \begin{itemize}
        \item Weniger gekoppelter Code
        \item Mehr Vertrauen in Änderungen
        \item Regressionen werden vermieden
    \end{itemize}
\end{itemize}
\subsubsection*{Kontinuierliche Intergration (Continuous Integration, CI)}
\begin{itemize}
    \item Entwickler \fatsf{integrieren Code mehmals täglich}
    \item Nach jedem Commit wird das System \fatsf{automatisch gebaut \& getestet}
    \item \fatsf{Voraussetzungen:}
    \begin{itemize}
        \item Versionskontrolle (Git, SVN, ...)
        \item Automatisierte Builds
        \item Automatisierte Tests
    \end{itemize}
\end{itemize}
\subsubsection*{Refactoring - Sauberer Code (Clean Code)}
\begin{itemize}
    \item \fatsf{Verbessern der Code-Struktur, ohne das Verhalten zu ändern}
    \item Frage vor neuen Features: Kann ich den Code zuerst verbessern?
    \item Falls ja $\rightarrow$ \fatsf{Refactoring durchführen, dann Feature hinzufügen}
\end{itemize}
\subsubsection*{Paarprogrammierung (Pair Programming)}
Zwei Entwickler arbeiten zusammen:
\begin{itemize}
    \item Einer schreibt den Code
    \item Der andere denkt strategisch mit
    \item Rollen wechseln regelmäßig $\rightarrow$ Verhindert Wissenssilos
\end{itemize}
Voraussetzung: \fatsf{Ähnliches Skill-Level der Entwickler}
\subsubsection*{Kollektive Code-Eigentümerschaft (Collective Ownership)}
\begin{itemize}
    \item \fatsf{Jeder darf jeden Teil des Codes ändern}
    \item Das ganze Team ist für die Code-Qualität verantwortlich
    \item Niemand wird für Fehler allein verantwortlich gemacht
\end{itemize}
\subsubsection*{Coding Standards}
\begin{itemize}
    \item Einheitliche Code-Richtilinien für das Team
    \item Wichtige Prinzipien:
    \begin{itemize}
        \item \fatsf{\endquote{No Duplicate Code}} $\rightarrow$ Keine Code-Dopplungen
        \item \fatsf{Einfachheit \& Verständlichkeit}
        \item Fördert Zusammenarbeit (Pair Programming, Refactoring)
    \end{itemize}
\end{itemize}
\subsubsection*{Fazit - Wann sollte man XP nutzen?}
\begin{itemize}
    \item \fatsf{Gut geeignet für:}
    \begin{itemize}
        \item Kleine bis mittelgroße Teams (3-10 Personen)
        \item Dynamische Prjekte mit \fatsf{häufigen Anforderungsänderungen}
        \item Teams mit \fatsf{hoher Disziplin \& enger Kundenbeteiligung}
    \end{itemize}
    \item \fatsf{Nicht geeignet für:}
    \begin{itemize}
        \item Große, strak regulierte Projekte (z.B. Medizintechnik, Luftfahrt)
        \item Projekte mit \fatsf{festen, unveränderlichen Anforderungen}
        \item Teams mit \fatsf{wenig Kundenkommunikation}
    \end{itemize}
    \fatsf{Fazit:} XP fördert schnelle, qualitativ hochwertige Entwicklung - erfordert aber eine enge Zusammenarbeit und disziplinierte Prozesse
\end{itemize}
\subsection{Part V: Software Engineering: Closing Remarks}
\subsubsection*{Fehlende Intuition bei Software}
\begin{itemize}
    \item Software ist weniger intuitiv als physische Materialien
    \item Kunden, Nutzer und sogar Entwickler haben oft \fatsf{falsche Vorstellungen} über die Eigenschaften von Software
    \item Konsequenz: \fatsf{Missverständnisse und Fehleinschätzungen}
\end{itemize}
\subsubsection*{Die Malleabilitätsillusion (\enquote{Software ist formbar})}
\begin{itemize}
    \item \fatsf{Software wirkt flexibel, ist es aber nicht!}
    \item Sobald Software \fatsf{einmal existiert und Abhängigkeiten hat}, wird sie \fatsf{so starr wie Stein}
    \item Änderungen sind dann \fatsf{sehr teuer \& aufwendig}
    \item \fatsf{Lösung:}
    \begin{itemize}
        \item Sorgfältige Planung \& Design
        \item Wartungsfreundliche Prozesse nutzen
        \item Umfassende Dokumentation sicherstellen
    \end{itemize}
\end{itemize}
\subsubsection*{Form $\neq$ Funktion in der Softwareentwicklung}
\begin{itemize}
    \item Bei physischen Objekten gibt die Form oft Hinweise auf die Funktion
    \item Bei Software ist das \fatsf{nicht} der Fall:
    \begin{itemize}
        \item Verschiedene Code-Darstellungen einer Funktion sehen unterschiedlich aus, \fatsf{machen aber dasselbe}
    \end{itemize}
    \item \fatsf{Warum ist die Funktion von Software nicht aus ihrem Code ablesbar?}
    \begin{enumerate}
        \item Code ist \fatsf{abstrakt \& textbasiert}
        \item Die Funktion wird oft erst \fatsf{während der Ausführung sichtbar}
    \end{enumerate}
    \item \fatsf{Konsequenzen:}
    \begin{itemize}
        \item Spezifikation \& Dokumentation sind essentiell (wichtiger als in traditionellen Ingenieurwissenschaften)
        \item Grafische Notation (z.B. UML können helfen)
        \item Testgetriebene Entwicklung (TDD) hilft, Funktionalität aus Anwendungsfällen abzuleiten
        \item Service-orientierte Architekturen (SOA) können Architektur und Funktionalität besser in Einklang bringen
    \end{itemize}
\end{itemize}
\subsubsection*{Wichtige Erkenntnis für Software Engineering}
\begin{itemize}
    \item Software-Entwicklung erfordert einen \fatsf{ingenieurwissenschaftlichen Ansatz}
    \item Muss auf \fatsf{systematischen, wissenschaftlichen Prinzipien} basieren
    \item Ist \fatsf{alternativlos}, weil Software ein \fatsf{immaterielles \& abstraktes Konstruktionsmaterial} ist
\end{itemize}
\fatsf{Fazit:} Softwareentwicklung ist kein reines \enquote{Handwerk}, sondern eine \fatsf{ingenieurmäßige Disziplin}, die durch präzise Planung, Methoden und Prinzipien gesteuert werden muss